\documentclass[11pt]{article}
\usepackage[margin=1in]{geometry}
\usepackage[T1]{fontenc}
\usepackage[utf8]{inputenc}
\usepackage{amsmath,amssymb,amsthm,mathtools}
\usepackage{booktabs}
\usepackage{graphicx}
\usepackage{xcolor}
\usepackage{microtype}
\usepackage[colorlinks=true,linkcolor=blue!55!black,citecolor=blue!55!black,urlcolor=blue!55!black]{hyperref}

\newtheorem{theorem}{Theorem}
\newtheorem{lemma}{Lemma}
\newcommand{\D}{\mathbb D}
\newcommand{\CM}{\mathrm{CM}}
\newcommand{\cH}{\mathcal H}
\newcommand{\Lap}{\mathcal L}

\title{A complete Mittag--Leffler classification for \texorpdfstring{$\mathcal H(K_p+K_q)$}{H(Kp+Kq)}}
\author{Candidate full-solution packet for arXiv:1703.02792, Question 3.10(ii)}
\date{11 August 2026}

\begin{document}
\maketitle

\begin{center}
\fbox{\parbox{0.92\textwidth}{\textbf{Status: candidate full solution, likely valid, subject to human review.}
The operator question is classified exactly by the standard zero-free phase
boundary for the two-parameter Mittag--Leffler function.}}
\end{center}

\section{Source question and answer}

For $a>0$, Anand and Chavan define
\[
 K_a(z,w)=\sum_{k\geq0}(k+1)^a(z\bar w)^k,\qquad z,w\in\D.
\]
Question 3.10(ii) of \emph{On sum of two subnormal kernels},
arXiv:1703.02792, asks for a necessary and sufficient condition for
$M_z$ on $\cH(K_p+K_q)$ to be subnormal \cite{source}.  Their paper proves
subnormality when $0<p\leq q\leq p+1$, and proves for a gap of exactly two
that $\cH(K_p+K_{p+2})$ is subnormal exactly when $p\geq1$.

Write
\[
 m=\min\{p,q\},\qquad M=\max\{p,q\},\qquad r=M-m.
\]
Let $E_{\alpha,\beta}$ denote the two-parameter Mittag--Leffler function
\[
 E_{\alpha,\beta}(z)=\sum_{j=0}^{\infty}
 \frac{z^j}{\Gamma(\alpha j+\beta)}.
\]
For $1<r\leq2$, let $\phi(r)$ be the zero-free phase boundary on the
negative real axis: it is the least $\beta$ such that
\[
 E_{r,\beta'}(-x)\geq0
 \quad\text{for every }x\geq0\text{ and every }\beta'\geq\beta.
\tag{1}\label{eq:phi}
\]
This is the standard $\alpha$--$\beta$ boundary tabulated by Hanneken,
Achar, and Vaught \cite{phase}; a recent primary-source computation
reproduces the boundary and its values \cite{honain}.  In particular,
$\phi(2)=3$.

\begin{theorem}[Complete classification]\label{thm:classification}
For $p,q>0$, the multiplication operator $M_z$ on $\cH(K_p+K_q)$ is
subnormal precisely in the following cases:
\begin{enumerate}
\item $r=0$;
\item $0<r\leq1$;
\item $1<r\leq2$ and $M\geq\phi(r)$.
\end{enumerate}
It is never subnormal when $r>2$.
\end{theorem}

Thus the answer is elementary outside the transition strip $1<r<2$, and in
that strip the exact threshold is the already established Mittag--Leffler
zero-free curve.  Some representative values are
\[
\begin{array}{c|rrrrrr}
r&1.05&1.20&1.50&1.80&1.95&2.00\\ \hline
\phi(r)&1.05924&1.26674&1.79365&2.46779&2.86070&3.00000\\
\phi(r)-r&0.00924&0.06674&0.29365&0.66779&0.91070&1.00000.
\end{array}
\]
The last row is the cutoff for the smaller exponent $m$.  At $r=2$, the
condition $M\geq3$ is exactly $m\geq1$, recovering the source theorem.

\section{The discrete-to-continuous rigidity lemma}

The source reduces subnormality to the Hausdorff moment property of
\begin{equation}\label{eq:sequence}
 a_k=\frac{1}{(k+1)^m+(k+1)^M},\qquad k\geq0.
\end{equation}
The central point is that, for this particular sequence, sampling only at the
positive integers loses no information.

\begin{lemma}[Rigidity of the Hausdorff interpolation]\label{lem:rigidity}
Put
\[
 f(s)=\frac{1}{s^m+s^M}=\frac{s^{-m}}{1+s^r}\qquad(s>0),
\]
using the principal powers.  Then \eqref{eq:sequence} is a Hausdorff moment
sequence if and only if the inverse Laplace density
\begin{equation}\label{eq:density}
 h_{m,M}(t)=t^{M-1}E_{r,M}(-t^r),\qquad t>0,
\end{equation}
is nonnegative.
\end{lemma}

\begin{proof}
The standard transform identity gives
\begin{equation}\label{eq:laplace}
 \Lap\{t^{M-1}E_{r,M}(-t^r)\}(s)
 =\frac{s^{r-M}}{1+s^r}
 =\frac{s^{-m}}{1+s^r}=f(s).
\end{equation}
If $h_{m,M}\geq0$, then
\[
 a_k=f(k+1)=\int_0^\infty e^{-(k+1)t}h_{m,M}(t)\,dt
 =\int_0^1x^k h_{m,M}(-\log x)\,dx,
\]
so $(a_k)$ is a Hausdorff moment sequence.

Conversely, suppose $a_k=\int_{[0,1]}x^k\,d\mu(x)$ for a finite positive
measure $\mu$.  On the right half-plane put
\[
 F(z)=\int_{[0,1]}x^z\,d\mu(x),\qquad
 G(z)=f(z+1).
\]
Both are bounded and analytic for $\Re z>0$.  For $G$, note that every
principal-branch solution of $s^r=-1$ has modulus one and real part strictly
less than one, whereas $s=z+1$ has real part greater than one.  The function
$F-G$ vanishes at every positive integer.  A nonzero bounded analytic
function on a half-plane has zeros satisfying the Blaschke condition
\[
 \sum_n\frac{\Re z_n}{1+|z_n|^2}<\infty;
\]
the positive integers violate it.  Hence $F=G$ identically.

Push $\mu|_{(0,1]}$ forward under $x=e^{-t}$.  The equality $F=G$, together
with \eqref{eq:laplace}, says that its Laplace transform is also the transform
of the finite signed measure
\[
 e^{-t}h_{m,M}(t)\,dt.
\]
This signed measure is finite: the poles of $f$ all have real part strictly
less than one, and the branch-cut contribution has only algebraic growth.
Uniqueness of Laplace transforms of finite measures therefore identifies the
two measures.  Since the first is positive, $h_{m,M}\geq0$ almost everywhere,
and continuity upgrades this to pointwise nonnegativity.
\end{proof}

The bounded-analytic step is essential.  Complete monotonicity of the
continuous function $f$ would certainly suffice, but without the rigidity
lemma it would not automatically be necessary for an integer-sampled
Hausdorff sequence.

\section{Proof of the classification}

\begin{proof}[Proof of Theorem~\ref{thm:classification}]
If $r=0$, then
\[
 a_k=\frac{1}{2(k+1)^m}
 =\frac1{2\Gamma(m)}\int_0^1x^k(-\log x)^{m-1}\,dx,
\]
so $M_z$ is subnormal.

Assume first that $0<r\leq1$.  Since $M=m+r>r$, Schneider's sharp theorem
says that $E_{r,M}(-x)$ is completely monotone, hence nonnegative, on
$[0,\infty)$ \cite{schneider,simon}.  Lemma~\ref{lem:rigidity} gives
subnormality for every $m>0$.  This recovers and slightly reorganizes the
source's gap-at-most-one region.

Now let $1<r\leq2$.  By Lemma~\ref{lem:rigidity}, subnormality is equivalent
to
\[
 E_{r,M}(-x)\geq0\qquad(x\geq0).
\]
The positivity set is an upper interval in the second parameter.  Indeed,
for $\delta>0$ the fractional-integration identity is
\[
 t^{M+\delta-1}E_{r,M+\delta}(-t^r)
 =\frac1{\Gamma(\delta)}\int_0^t
 (t-u)^{\delta-1}u^{M-1}E_{r,M}(-u^r)\,du.
\]
Consequently, the exact threshold is the zero-free boundary $M=\phi(r)$
defined in \eqref{eq:phi}.  This proves the third case, including its
endpoint.  The identity $\phi(2)=3$ gives precisely the source's result for
$K_m+K_{m+2}$.

Finally suppose $r>2$.  The principal-branch denominator $1+s^r$ has a
dominant conjugate pair of simple poles
\[
 \zeta_\pm=e^{\pm i\pi/r},\qquad
 \Re\zeta_\pm=\cos(\pi/r)>0.
\]
Moving the inverse-Laplace contour left and collecting residues yields
\begin{equation}\label{eq:asymptotic}
 h_{m,M}(t)=\frac2r e^{t\cos(\pi/r)}
 \cos\!\left(t\sin(\pi/r)+\frac{(1-M)\pi}{r}\right)
 +o\!\left(e^{t\cos(\pi/r)}\right).
\end{equation}
All other poles have smaller real part, and the branch-cut term is
subexponential.  The leading cosine changes sign infinitely often.  Thus
$h_{m,M}$ is not nonnegative, and Lemma~\ref{lem:rigidity} rules out
subnormality for every $m>0$.
\end{proof}

\section{Checks, novelty, and limitations}

The bundled script evaluates the defining Mittag--Leffler series at high
precision.  It confirms sign changes just below representative tabulated
phase thresholds, positivity at sampled points just above them, the identity
$h_{1,3}(t)=1-\cos t$ at $r=2$, and the oscillatory obstruction for $r>2$.
These computations are sanity checks; the proof rests on the transform,
Blaschke, phase-boundary, and residue arguments.

The initially plausible elementary cutoff $m\geq r-1$ is not sharp.  For
example, at $r=1.5$ the true cutoff is
$m\geq\phi(1.5)-1.5\approx0.29365$, whereas $r-1=0.5$.  The phase-boundary
formulation is therefore mathematically substantive, not cosmetic.

On 11 August 2026, the run registry and cheap indexes were searched for
arXiv:1703.02792, the exact kernel notation, and the moment sequence
\eqref{eq:sequence}.  Bounded primary-source searches for the title, exact
question, and close combinations of ``subnormal kernel'', ``Hausdorff
moment'', and ``Mittag--Leffler'' found the source paper and the independent
special-function phase literature, but no later paper making this
operator-theoretic connection or answering Question 3.10(ii).  The novelty is
the exact reduction and synthesis; the special-function phase curve itself
is not claimed as new.

Human review should focus on (i) boundedness and zero-set uniqueness in
Lemma~\ref{lem:rigidity}; (ii) uniqueness of the signed inverse-Laplace
measure after the unit shift; and (iii) the interpretation and endpoint of
the published zero-free boundary $\phi$.  The result answers Question 3.10(ii)
only; Questions 3.10(i) and 3.11 remain outside its scope.

\appendix
\section{Source-page evidence}

Figure~\ref{fig:source} reproduces printed page 12 of the source PDF,
including Question 3.10(ii).

\begin{figure}[p]
\centering
\includegraphics[height=0.88\textheight]{figures/open_problem_crop.png}
\caption{Source paper, printed page 12: Question 3.10, including item (ii).}\label{fig:source}
\end{figure}

\begin{thebibliography}{9}
\bibitem{source}
A. Anand and S. Chavan,
\emph{On sum of two subnormal kernels},
arXiv:1703.02792v2 [math.FA], 2017.

\bibitem{phase}
J. W. Hanneken, B. N. Achar, and D. M. Vaught,
\emph{An alpha--beta phase diagram representation of the zeros and
properties of the Mittag--Leffler function},
Advances in Mathematical Physics (2013), Article ID 421685.

\bibitem{honain}
A. H. Honain, K. M. Furati, I. O. Sarumi, and A. Q. M. Khaliq,
\emph{Rational approximations for oscillatory two-parameter Mittag--Leffler
function}, arXiv:2312.07444, 2023; Fractal and Fractional 8 (2024), 319.

\bibitem{schneider}
W. R. Schneider,
\emph{Completely monotone generalized Mittag--Leffler functions},
Expositiones Mathematicae 14 (1996), 3--16.

\bibitem{simon}
T. Simon,
\emph{Mittag--Leffler functions and complete monotonicity},
arXiv:1312.4513, 2013.
\end{thebibliography}

\end{document}
